\begin{claim}{B.15a}{Proved}
Put $F=k(\mu_\infty)$, where $k/\Q$ is finitely generated. For a
hyperbolic curve $Y/F$, an arbitrary continuous section of
$\pi_1(Y)\to G_F$ fixes at most one lifted geometric cusp. No extension
of this section to $G_k$ is required.
\end{claim}
\begin{proof}
We give the arithmetic input and the finite-cover argument explicitly.
For every finite extension $F'/F$ and every abelian variety $A/F'$, one has
\[
 T(A)^{G_{F'}}=0,\qquad \bigcap_{n\ge1}nA(F')=0.
\]
Here is the spreading argument extending Ribet's number-field theorem
\citep[Appendix, Theorem 1]{BC_KatzLang}. Write
$F'=k_0(\mu_\infty)$ with $k_0/\Q$ finitely generated and descend the
finite data to $k_0$. Let $a$ be its algebraic constant number field, and
spread $A$ to an abelian scheme over a smooth geometrically integral
$a$-variety $S$ with function field $k_0$. A rational torsion point over
$F'$ extends over $S_{a(\mu_\infty)}$: its finite etale division scheme
has a generic section, and normality extends that section. Evaluation at
a closed point of $S$, after passage to its number-field residue and its
cyclotomic closure, is injective on torsion sections: the equalizer of two
sections of a finite etale scheme is open and closed. Ribet's finiteness
therefore gives the first assertion. This argument also works after every
finite extension of $F$.

For the second assertion first consider a number-field cyclotomic closure
$a(\mu_\infty)$. A point divisible by every integer has compatible
factorial division points, by compactness of their finite nonempty sets.
Descend the abelian variety and the point to a finite subextension $a_1$.
The Kummer cocycle on $G_{a_1}$ vanishes on the normal subgroup
$G_{a(\mu_\infty)}$. Its values are therefore fixed by that subgroup,
so the first assertion makes the cocycle zero. All division points are
then defined over $a_1$. Reduction in a completion at a good finite place,
modulo successive powers of its maximal ideal, puts the point in the
intersection of the kernels of finite quotient groups. It is zero.
For general $F'$, spread a proposed nonzero infinitely divisible point as
a section over $S$ and remove its intersection with the identity section.
Every division point extends over $S_{a(\mu_\infty)}$ by the same finite
etale argument. Specialization gives a nonzero infinitely divisible point
over a number-field cyclotomic closure, a contradiction.

Suppose now that two distinct lifted cusps are fixed by the section
$\sigma:G_F\to\pi_1(Y)$. A finite geometric etale cover separates their
images. Refine it by a characteristic cover whose proper compactification
has genus at least two. Such a geometric cover exists: for a punctured
genus-zero curve, choose three punctures and restrict a Fermat quartic
cover of the resulting three-punctured line;
for genus one, an etale cover first supplies two boundary points and
a double cover branched there has genus two; in proper genus at least two
the identity suffices. A characteristic refinement of the Galois closure
preserves this genus bound by Riemann--Hurwitz. The subgroup generated by
this characteristic geometric subgroup and $\sigma(G_F)$ descends the
cover to $F$. The two fixed cusps become distinct $F$-rational points
$a,b$ of its proper compactification $\overline Y$.

Both points induce the same section of $\pi_1(\overline Y)\to G_F$,
because cusp inertia is killed on compactification. Under the Abel map
$z\mapsto[z-b]$, their difference has zero Kummer class for every
multiplication cover of $\Jac(\overline Y)$, hence
$[a-b]\in\bigcap_n n\Jac(\overline Y)(F)=0$. A function with divisor
$a-b$ would give a degree-one map to $\mathbb P^1$, contrary to the positive
genus. This proves uniqueness, including distinct lifts of the same base
cusp. It is the finite-neighbourhood proof of Stix's packet theorem,
with its arithmetic hypothesis verified over $F$
\citep[Theorem 14 and its proof; Appendix]{CF_Stix2008}.
\end{proof}

\begin{claim}{B.15u}{Proved}
Let $(\mathcal C,\mathcal D)/R$ be a split stable marked model in the
characteristic-zero setting of Part B. Let $x\in U(K)$ and let
$s:G_k\to G_K$ be any residue splitting. For the genuine action specified
by the point section, put
$H=(s_x\circ s)(G_k^\circ)$. Then
\[
 C_\Pi(H)=
 \begin{cases}
  1,&\text{if $x$ specializes to a smooth unmarked point},\\
  I_c,&\text{if $x$ specializes to a marked point $c$},\\
  \Pi_e,&\text{if $x$ specializes to a node $e$}.
 \end{cases}
\]
The inertia groups are the actual lifted subgroups selected by the local
point path. The equality is independent of the positive contact
multiplicity at a mark and of the two positive contact orders
$a,m_e-a$ at a node. In particular no other lifted node group or cusp
group is contained in this centralizer.
\end{claim}
\begin{proof}
We specify both the arithmetic vertex groups and the profinite-tree
argument; passing to an abstract tree or assuming that arbitrary fixed
subgraphs are path connected would not suffice.

Remove the open edges (the marked cusps) from the geometric semi-graph of
B.3, retaining the vertex groups and the closed-edge groups. This finite
graph of profinite groups has injective vertex and edge maps. Its standard
Bass--Serre graph $\mathcal T$ is a profinite tree, with
\[
 V(\mathcal T)=\coprod_v\Pi/\Pi_v,
 \qquad E(\mathcal T)=\coprod_e\Pi/\Pi_e.
\]
The sums are finite, so the edge set is compact. The geometric stabilizers
are the indicated conjugates. These are the construction and exact tree
theorem of \citet[\S3.7 and Proposition 3.8, pp.~419--420]{BC_ZM1989}.
The one-vertex graph with no nodes is included directly. The split
arithmetic action acts continuously on this tree and preserves the two
branches of each lifted edge; in particular there are no edge inversions.

For a lifted vertex $v$, let $Y_v$ be its normalized component minus the
nodes and marks. Its vertical logarithmic chart is the constant chart
$\mathbb N\to\mathcal O_{Y_v}$, $1\mapsto0$, pulled back from the base
log point. The Kummer covers of this chart are the roots of the base
parameter. After adjoining these roots the remaining finite cover is
etale on $Y_v$. Thus the arithmetic logarithmic vertex group $D_v$ has
the Cartesian description
\[
 D_v\simeq\pi_1(Y_v)\times_{G_k}G_K.
 \tag{B.vertex}
\]
The projection to $\pi_1(Y_v)$ kills precisely the vertical inertia,
which intersects the geometric vertex group trivially and centralizes it.
This is also the vertex part of the normalized-component gluing in B.3:
the chart is strict over the base log point away from its deleted special
points, and adjoining a root of the base chart leaves an ordinary finite
etale component cover. Conversely these component covers and the compatible
base roots give every such local Kummer cover. The injections in B.3
identify its geometric kernel with $\Pi_v$, and its projection onto $G_K$
is surjective. Its image is the entire arithmetic stabilizer: any other
stabilizer element can be multiplied by a local element with the same
image in $G_K$ to lie in $N_\Pi(\Pi_v)=\Pi_v$, the terminality assertion
of \citet[Proposition 1.2]{BC_MochComb}. Conjugating the chosen lift
conjugates this entire Cartesian description.

Pulling (B.vertex) back along $s$ identifies
$D_v\cap E_s$ with $\pi_1(Y_v)$. Consequently, for \emph{every}
$H$-fixed vertex, even one not fixed by the whole point-section image of
$G_k$, projection gives an actual continuous section
\[
 G_{k(\mu_\infty)}\longrightarrow
 \pi_1((Y_v)_{k(\mu_\infty)}).
\]
It projects to the identity on $G_{k(\mu_\infty)}$ because $H$ already
does so. The incident edges at $v$ are the lifted cusps of $Y_v$
corresponding to node branches; the identification follows from their
coset descriptions $\Pi_v/\Pi_e$. Distinct incident edges give distinct
lifted cusps, including when their images are the same branch at a finite
level. This identification is equivariant for $D_v$. Therefore two
distinct incident $H$-fixed edges would give two fixed lifted cusps for
the displayed section, contrary to B.15a. Every vertex of
$\mathcal T^H$ has degree at most one.

Theorem 2.8 of \citet[pp.~413--414]{BC_ZM1989} applies to an arbitrary
profinite group acting on a full profinite tree; it gives that
$\mathcal T^H$ is a subtree. It is nonempty: the local point path fixes
the vertex of the specialization if it is smooth, and fixes the lifted
node edge if the specialization is nodal. Its edge set is closed in the
compact edge set of $\mathcal T$, hence compact. A connected profinite
graph with compact edges and more than one vertex has an edge incident to
every vertex \citep[Lemma 2.14, p.~417]{BC_ZM1989}.
For clarity, this last fact follows by taking finite graph quotients:
the closed subsets of edges whose endpoint agrees with a given vertex
in each quotient have the finite intersection property.

It follows that $\mathcal T^H$ is either one vertex or one edge. Indeed,
in the nontrivial case every vertex has degree exactly one. The two endpoint
maps identify its vertex set with two disjoint copies of its compact edge
space: they are continuous bijections onto disjoint closed sets covering
all vertices, hence homeomorphisms. If the edge space had two points,
a nontrivial clopen partition would give a disconnected finite quotient
graph. This contradicts connectedness. This argument does not collapse
an edge and does not presume compactness of the edges after such a collapse.

Write $C=C_\Pi(H)$. It preserves $\mathcal T^H$, since its elements
commute with $H$. At a node the fixed tree already contains the local
edge, so it is exactly that edge. The action has no inversions, hence
$C$ lies in its geometric stabilizer $\Pi_e$. B.15 gives the reverse
inclusion for every splitting and for all contact orders.

If the specialization is smooth and unmarked, the action at its fixed
vertex, under (B.vertex), is the ordinary point action of $Y_v$ at that
special point. R15 gives $C\cap\Pi_v=1$. An incident fixed edge would
have its nontrivial inertia centralized by $H$: the edge action is the
cyclotomic character, which is trivial on $G_k^\circ$. Thus no such edge
exists, and the fixed tree is just $v$. Since $C$ preserves it,
$C\subset\Pi_v$ and $C=1$.

If the specialization is the mark $c$, the action at the fixed vertex
is the cuspidal action computed in B.13. R15 gives
$C\cap\Pi_v=I_c$. An incident fixed node edge would instead give a second
fixed lifted cusp of $Y_v$, contradicting B.15a. Again the fixed tree is
just $v$, and consequently $C=I_c$. Distinct lifted peripheral groups
are distinct maximal procyclic groups by the peripheral terminality in
\citet[Proposition 1.2]{BC_MochComb}; hence the displayed equalities
exclude every other node or cusp group. The multiplicities affect the
local Kummer homomorphism, but B.15 and the tree upper bounds show that
they do not affect this centralizer.
\end{proof}
